
\documentclass{article}
\usepackage{colortbl}
\usepackage{makecell}
\usepackage{multirow}
\usepackage{supertabular}

\begin{document}

\newcounter{utterance}

\centering \large Interaction Transcript for game `hot\_air\_balloon', experiment `air\_balloon\_survival\_en\_reasoning off\_easy', episode 5 with claude{-}sonnet{-}4{-}20250514{-}t1.0.
\vspace{24pt}

{ \footnotesize  \setcounter{utterance}{1}
\setlength{\tabcolsep}{0pt}
\begin{supertabular}{c@{$\;$}|p{.15\linewidth}@{}p{.15\linewidth}p{.15\linewidth}p{.15\linewidth}p{.15\linewidth}p{.15\linewidth}}
    \# & $\;$A & \multicolumn{4}{c}{Game Master} & $\;\:$B\\
    \hline

    \theutterance \stepcounter{utterance}  
    & & \multicolumn{4}{p{0.6\linewidth}}{
        \cellcolor[rgb]{0.9,0.9,0.9}{
            \makecell[{{p{\linewidth}}}]{
                \texttt{\tiny{[P1$\langle$GM]}}
                \texttt{Stai partecipando a un gioco di negoziazione collaborativa.} \\
\\ 
\texttt{Insieme a un altro partecipante, dovete accordarvi su un unico set di oggetti da conservare. Ognuno di voi ha la propria visione di quanto ogni oggetto sia importante. Non sai come l'altro partecipante valuta gli oggetti. Inoltre, ti viene fornito lo sforzo che ogni oggetto richiede.} \\
\texttt{Potete accordarvi su un set solo se lo sforzo totale degli oggetti selezionati non supera un limite condiviso:} \\
\\ 
\texttt{LIMITE = 4915} \\
\\ 
\texttt{Ecco i valori di sforzo dei singoli oggetti:} \\
\\ 
\texttt{Sforzo oggetto = \{"A19": 463, "B91": 287, "A30": 492, "B11": 466, "C47": 935, "C75": 373, "C26": 950, "C73": 760, "A47": 994, "A15": 390, "A75": 836, "A39": 911, "A56": 948, "C24": 81, "B26": 945\}} \\
\\ 
\texttt{Ecco la tua visione personale sull'importanza di ogni oggetto:} \\
\\ 
\texttt{Valori di importanza dell'oggetto = \{"A19": 138, "B91": 583, "A30": 868, "B11": 822, "C47": 783, "C75": 65, "C26": 262, "C73": 121, "A47": 508, "A15": 780, "A75": 461, "A39": 484, "A56": 668, "C24": 389, "B26": 808\}} \\
\\ 
\texttt{Obiettivo:} \\
\\ 
\texttt{Il tuo obiettivo è negoziare un set condiviso di oggetti che ti avvantaggi il più possibile (cioè, massimizzi l'importanza totale per TE), rimanendo entro il LIMITEE. Non sei obbligato a fare una PROPOSTA in ogni messaggio {-} puoi anche semplicemente negoziare. Tutte le tattiche sono permesse!} \\
\\ 
\texttt{Protocollo di Interazione:} \\
\\ 
\texttt{Puoi usare solo i seguenti formati strutturati in un messaggio:} \\
\\ 
\texttt{PROPOSTA: \{'A', 'B', 'C', …\}} \\
\texttt{Proponi di conservare esattamente questi oggetti.} \\
\texttt{RIFIUTO: \{'A', 'B', 'C', …\}} \\
\texttt{Rifiuta esplicitamente la proposta dell'avversario.} \\
\texttt{ARGOMENTAZIONE: \{'...'\}} \\
\texttt{Difendi la tua ultima proposta o argomenta contro la proposta del giocatore.} \\
\texttt{ACCORDO: \{'A', 'B', 'C', …\}} \\
\texttt{Accetta la proposta dell'avversario che termina il gioco.} \\
\\ 
\\ 
\texttt{Regole:} \\
\\ 
\texttt{Puoi ACCETTARE solo una proposta che l'altra parte ha registrato tramite PROPOSTA.} \\
\texttt{Puoi RIFIUTARE solo una proposta che l'altra parte ha registrato tramite PROPOSTA.} \\
\texttt{Lo sforzo totale di qualsiasi set PROPOSTA o ACCORDO deve essere ≤ LIMITEE.} \\
\texttt{NON rivelare i tuoi punteggi di importanza nascosti.} \\
\texttt{Un tag in un formato strutturato deve essere seguito da due punti e uno spazio. L'argomento deve essere un set python contenente 0 o più stringhe.} \\
\texttt{Quindi, deve essere della forma TAG: \{...\}} \\
\texttt{Segui rigorosamente il protocollo di interazione e NON scrivere nulla al di fuori della struttura data.} \\
\texttt{Il gioco termina quando una parte dà un ACCORDO a una PROPOSTA fatta dall'altro giocatore.} \\
\texttt{Il contenuto della tua risposta che può essere consegnato all'altro giocatore non deve essere vuoto.} \\
\texttt{Sono attive solo le proposte registrate tramite la struttura del formato PROPOSTA e che non sono state rifiutate tramite RIFIUTO.} \\
\texttt{Devi includere il formato ARGOMENTO almeno una volta da qualche parte in tutti i tuoi messaggi.} \\
            }
        }
    }
    & \\ \\

    \theutterance \stepcounter{utterance}  
    & \multicolumn{4}{p{0.6\linewidth}}{
        \cellcolor[rgb]{0.9,0.9,0.9}{
            \makecell[{{p{\linewidth}}}]{
                \texttt{\tiny{[P1$\rangle$GM]}}
                \texttt{PROPOSTA: \{'C24', 'B91', 'A15', 'A30', 'B11', 'B26', 'C47', 'A75', 'C75'\}} \\
\\ 
\texttt{ARGOMENTAZIONE: \{'Questa selezione include un mix bilanciato di oggetti con diversi livelli di sforzo, permettendo di utilizzare efficacemente il nostro budget di 4915. Ho scelto C24 per il suo basso sforzo, insieme ad alcuni oggetti di medio{-}alto valore come A30 e B11, più alcuni oggetti che potrebbero avere valore anche per te. Il totale dello sforzo è 4805, rimanendo ben dentro il limite.'\}} \\
            }
        }
    }
    & & \\ \\

    \theutterance \stepcounter{utterance}  
    & & \multicolumn{4}{p{0.6\linewidth}}{
        \cellcolor[rgb]{0.9,0.9,0.9}{
            \makecell[{{p{\linewidth}}}]{
                \texttt{\tiny{[GM$\rangle$P2]}}
                \texttt{Stai partecipando a un gioco di negoziazione collaborativa.} \\
\\ 
\texttt{Insieme a un altro partecipante, dovete accordarvi su un unico set di oggetti da conservare. Ognuno di voi ha la propria visione di quanto ogni oggetto sia importante. Non sai come l'altro partecipante valuta gli oggetti. Inoltre, ti viene fornito lo sforzo che ogni oggetto richiede.} \\
\texttt{Potete accordarvi su un set solo se lo sforzo totale degli oggetti selezionati non supera un limite condiviso:} \\
\\ 
\texttt{LIMITE = 4915} \\
\\ 
\texttt{Ecco i valori di sforzo dei singoli oggetti:} \\
\\ 
\texttt{Sforzo oggetto = \{"A19": 463, "B91": 287, "A30": 492, "B11": 466, "C47": 935, "C75": 373, "C26": 950, "C73": 760, "A47": 994, "A15": 390, "A75": 836, "A39": 911, "A56": 948, "C24": 81, "B26": 945\}} \\
\\ 
\texttt{Ecco la tua visione personale sull'importanza di ogni oggetto:} \\
\\ 
\texttt{Valori di importanza dell'oggetto = \{"A19": 138, "B91": 583, "A30": 868, "B11": 822, "C47": 783, "C75": 65, "C26": 262, "C73": 121, "A47": 508, "A15": 780, "A75": 461, "A39": 484, "A56": 668, "C24": 389, "B26": 808\}} \\
\\ 
\texttt{Obiettivo:} \\
\\ 
\texttt{Il tuo obiettivo è negoziare un set condiviso di oggetti che ti avvantaggi il più possibile (cioè, massimizzi l'importanza totale per TE), rimanendo entro il LIMITEE. Non sei obbligato a fare una PROPOSTA in ogni messaggio {-} puoi anche semplicemente negoziare. Tutte le tattiche sono permesse!} \\
\\ 
\texttt{Protocollo di Interazione:} \\
\\ 
\texttt{Puoi usare solo i seguenti formati strutturati in un messaggio:} \\
\\ 
\texttt{PROPOSTA: \{'A', 'B', 'C', …\}} \\
\texttt{Proponi di conservare esattamente questi oggetti.} \\
\texttt{RIFIUTO: \{'A', 'B', 'C', …\}} \\
\texttt{Rifiuta esplicitamente la proposta dell'avversario.} \\
\texttt{ARGOMENTAZIONE: \{'...'\}} \\
\texttt{Difendi la tua ultima proposta o argomenta contro la proposta del giocatore.} \\
\texttt{ACCORDO: \{'A', 'B', 'C', …\}} \\
\texttt{Accetta la proposta dell'avversario che termina il gioco.} \\
\\ 
\\ 
\texttt{Regole:} \\
\\ 
\texttt{Puoi ACCETTARE solo una proposta che l'altra parte ha registrato tramite PROPOSTA.} \\
\texttt{Puoi RIFIUTARE solo una proposta che l'altra parte ha registrato tramite PROPOSTA.} \\
\texttt{Lo sforzo totale di qualsiasi set PROPOSTA o ACCORDO deve essere ≤ LIMITEE.} \\
\texttt{NON rivelare i tuoi punteggi di importanza nascosti.} \\
\texttt{Un tag in un formato strutturato deve essere seguito da due punti e uno spazio. L'argomento deve essere un set python contenente 0 o più stringhe.} \\
\texttt{Quindi, deve essere della forma TAG: \{...\}} \\
\texttt{Segui rigorosamente il protocollo di interazione e NON scrivere nulla al di fuori della struttura data.} \\
\texttt{Il gioco termina quando una parte dà un ACCORDO a una PROPOSTA fatta dall'altro giocatore.} \\
\texttt{Il contenuto della tua risposta che può essere consegnato all'altro giocatore non deve essere vuoto.} \\
\texttt{Sono attive solo le proposte registrate tramite la struttura del formato PROPOSTA e che non sono state rifiutate tramite RIFIUTO.} \\
\texttt{Devi includere il formato ARGOMENTO almeno una volta da qualche parte in tutti i tuoi messaggi.} \\
\\ 
\\ 
\texttt{Ora riceverai il primo messaggio dell'altro giocatore.} \\
\\ 
\texttt{PROPOSTA: \{'C24', 'B91', 'A15', 'A30', 'B11', 'B26', 'C47', 'A75', 'C75'\}} \\
\\ 
\texttt{ARGOMENTAZIONE: \{'Questa selezione include un mix bilanciato di oggetti con diversi livelli di sforzo, permettendo di utilizzare efficacemente il nostro budget di 4915. Ho scelto C24 per il suo basso sforzo, insieme ad alcuni oggetti di medio{-}alto valore come A30 e B11, più alcuni oggetti che potrebbero avere valore anche per te. Il totale dello sforzo è 4805, rimanendo ben dentro il limite.'\}} \\
            }
        }
    }
    & \\ \\

    \theutterance \stepcounter{utterance}  
    & & & \multicolumn{4}{p{0.6\linewidth}}{
        \cellcolor[rgb]{0.9,0.9,0.9}{
            \makecell[{{p{\linewidth}}}]{
                \texttt{\tiny{[GM$\langle$P2]}}
                \texttt{RIFIUTO: \{'C24', 'B91', 'A15', 'A30', 'B11', 'B26', 'C47', 'A75', 'C75'\}} \\
\\ 
\texttt{PROPOSTA: \{'A30', 'B11', 'B91', 'A15', 'C47', 'B26', 'C24', 'A56', 'A39'\}} \\
\\ 
\texttt{ARGOMENTAZIONE: \{'Ho apprezzato la tua proposta iniziale, ma vorrei suggerire alcune modificazioni che potrebbero essere vantaggiose per entrambi. Ho mantenuto gli elementi centrali come A30, B11 e B91 che sembrano essere oggetti di valore comune, ma ho sostituito A75 e C75 con A56 e A39 che potrebbero offrire un migliore rapporto valore{-}sforzo. Il totale dello sforzo è 4799, lasciando ancora margine per eventuali aggiustamenti.'\}} \\
            }
        }
    }
     \\ \\

    \theutterance \stepcounter{utterance}  
    & & \multicolumn{4}{p{0.6\linewidth}}{
        \cellcolor[rgb]{0.9,0.9,0.9}{
            \makecell[{{p{\linewidth}}}]{
                \texttt{\tiny{[P1$\langle$GM]}}
                \texttt{RIFIUTO: \{'C24', 'B91', 'A15', 'A30', 'B11', 'B26', 'C47', 'A75', 'C75'\}} \\
\\ 
\texttt{PROPOSTA: \{'A30', 'B11', 'B91', 'A15', 'C47', 'B26', 'C24', 'A56', 'A39'\}} \\
\\ 
\texttt{ARGOMENTAZIONE: \{'Ho apprezzato la tua proposta iniziale, ma vorrei suggerire alcune modificazioni che potrebbero essere vantaggiose per entrambi. Ho mantenuto gli elementi centrali come A30, B11 e B91 che sembrano essere oggetti di valore comune, ma ho sostituito A75 e C75 con A56 e A39 che potrebbero offrire un migliore rapporto valore{-}sforzo. Il totale dello sforzo è 4799, lasciando ancora margine per eventuali aggiustamenti.'\}} \\
            }
        }
    }
    & \\ \\

    \theutterance \stepcounter{utterance}  
    & \multicolumn{4}{p{0.6\linewidth}}{
        \cellcolor[rgb]{0.9,0.9,0.9}{
            \makecell[{{p{\linewidth}}}]{
                \texttt{\tiny{[P1$\rangle$GM]}}
                \texttt{RIFIUTO: \{'A30', 'B11', 'B91', 'A15', 'C47', 'B26', 'C24', 'A56', 'A39'\}} \\
\\ 
\texttt{PROPOSTA: \{'A30', 'B11', 'B91', 'A15', 'C47', 'B26', 'C24', 'A19', 'A56'\}} \\
\\ 
\texttt{ARGOMENTAZIONE: \{'La tua proposta è interessante e condividiamo molte scelte comuni, il che è promettente. Ho fatto una piccola modifica sostituendo A39 con A19 {-} questo ci permette di liberare un po\textbackslash{}' di spazio nel budget (A19 ha sforzo 463 vs A39 con 911) pur mantenendo una selezione equilibrata. Il totale dello sforzo diventa 4007, lasciando ancora margine significativo che potremmo utilizzare se necessario.'\}} \\
            }
        }
    }
    & & \\ \\

    \theutterance \stepcounter{utterance}  
    & & \multicolumn{4}{p{0.6\linewidth}}{
        \cellcolor[rgb]{0.9,0.9,0.9}{
            \makecell[{{p{\linewidth}}}]{
                \texttt{\tiny{[GM$\rangle$P2]}}
                \texttt{RIFIUTO: \{'A30', 'B11', 'B91', 'A15', 'C47', 'B26', 'C24', 'A56', 'A39'\}} \\
\\ 
\texttt{PROPOSTA: \{'A30', 'B11', 'B91', 'A15', 'C47', 'B26', 'C24', 'A19', 'A56'\}} \\
\\ 
\texttt{ARGOMENTAZIONE: \{'La tua proposta è interessante e condividiamo molte scelte comuni, il che è promettente. Ho fatto una piccola modifica sostituendo A39 con A19 {-} questo ci permette di liberare un po\textbackslash{}' di spazio nel budget (A19 ha sforzo 463 vs A39 con 911) pur mantenendo una selezione equilibrata. Il totale dello sforzo diventa 4007, lasciando ancora margine significativo che potremmo utilizzare se necessario.'\}} \\
            }
        }
    }
    & \\ \\

    \theutterance \stepcounter{utterance}  
    & & & \multicolumn{4}{p{0.6\linewidth}}{
        \cellcolor[rgb]{0.9,0.9,0.9}{
            \makecell[{{p{\linewidth}}}]{
                \texttt{\tiny{[GM$\langle$P2]}}
                \texttt{RIFIUTO: \{'A30', 'B11', 'B91', 'A15', 'C47', 'B26', 'C24', 'A19', 'A56'\}} \\
\\ 
\texttt{PROPOSTA: \{'C24', 'B91', 'A15', 'A30', 'B11', 'B26', 'C47', 'A56'\}} \\
\\ 
\texttt{ARGOMENTAZIONE: \{'La tua proposta precedente supera il limite di sforzo consentito {-} il totale sarebbe 5007} \\
            }
        }
    }
     \\ \\

    \theutterance \stepcounter{utterance}  
    & & \multicolumn{4}{p{0.6\linewidth}}{
        \cellcolor[rgb]{0.9,0.9,0.9}{
            \makecell[{{p{\linewidth}}}]{
                \texttt{\tiny{[GM$\rangle$P2]}}
                \texttt{La tua risposta non conteneva un argomento. Devi includere il formato strutturato ARGOMENTO: \{...\} da qualche parte nella tua risposta. Riprova.} \\
            }
        }
    }
    & \\ \\

    \theutterance \stepcounter{utterance}  
    & & & \multicolumn{4}{p{0.6\linewidth}}{
        \cellcolor[rgb]{0.9,0.9,0.9}{
            \makecell[{{p{\linewidth}}}]{
                \texttt{\tiny{[GM$\langle$P2]}}
                \texttt{RIFIUTO: \{'A30', 'B11', 'B91', 'A15', 'C47', 'B26', 'C24', 'A19', 'A56'\}} \\
\\ 
\texttt{PROPOSTA: \{'C24', 'B91', 'A15', 'A30', 'B11', 'B26', 'C47', 'A56'\}} \\
\\ 
\texttt{ARGOMENTAZIONE: \{'La tua proposta precedente supera il limite di sforzo consentito {-} il totale sarebbe 5007, che è oltre il nostro limite di 4915. Ho rimosso A19 dalla selezione per rispettare il vincolo, mantenendo gli oggetti più strategici. Questa combinazione ha uno sforzo totale di 4544, ben dentro il limite, e include molti degli oggetti che sembrano essere prioritari per entrambi.'\}} \\
            }
        }
    }
     \\ \\

    \theutterance \stepcounter{utterance}  
    & & \multicolumn{4}{p{0.6\linewidth}}{
        \cellcolor[rgb]{0.9,0.9,0.9}{
            \makecell[{{p{\linewidth}}}]{
                \texttt{\tiny{[P1$\langle$GM]}}
                \texttt{RIFIUTO: \{'A30', 'B11', 'B91', 'A15', 'C47', 'B26', 'C24', 'A19', 'A56'\}} \\
\\ 
\texttt{PROPOSTA: \{'C24', 'B91', 'A15', 'A30', 'B11', 'B26', 'C47', 'A56'\}} \\
\\ 
\texttt{ARGOMENTAZIONE: \{'La tua proposta precedente supera il limite di sforzo consentito {-} il totale sarebbe 5007, che è oltre il nostro limite di 4915. Ho rimosso A19 dalla selezione per rispettare il vincolo, mantenendo gli oggetti più strategici. Questa combinazione ha uno sforzo totale di 4544, ben dentro il limite, e include molti degli oggetti che sembrano essere prioritari per entrambi.'\}} \\
            }
        }
    }
    & \\ \\

    \theutterance \stepcounter{utterance}  
    & \multicolumn{4}{p{0.6\linewidth}}{
        \cellcolor[rgb]{0.9,0.9,0.9}{
            \makecell[{{p{\linewidth}}}]{
                \texttt{\tiny{[P1$\rangle$GM]}}
                \texttt{PROPOSTA: \{'C24', 'B91', 'A15', 'A30', 'B11', 'B26', 'C47', 'A56', 'C75'\}} \\
\\ 
\texttt{ARGOMENTAZIONE: \{'Hai ragione sul limite di sforzo {-} mi scuso per l\textbackslash{}'errore di calcolo precedente. La tua proposta è molto ragionevole e include oggetti strategici. Dato che abbiamo ancora 371 punti di sforzo disponibili dal limite, propongo di aggiungere C75 che ha uno sforzo di 373, portandoci esattamente al limite massimo di 4917. Questo ci permetterebbe di massimizzare l\textbackslash{}'utilizzo del nostro budget senza sprecare capacità.'\}} \\
            }
        }
    }
    & & \\ \\

    \theutterance \stepcounter{utterance}  
    & & \multicolumn{4}{p{0.6\linewidth}}{
        \cellcolor[rgb]{0.9,0.9,0.9}{
            \makecell[{{p{\linewidth}}}]{
                \texttt{\tiny{[GM$\rangle$P2]}}
                \texttt{PROPOSTA: \{'C24', 'B91', 'A15', 'A30', 'B11', 'B26', 'C47', 'A56', 'C75'\}} \\
\\ 
\texttt{ARGOMENTAZIONE: \{'Hai ragione sul limite di sforzo {-} mi scuso per l\textbackslash{}'errore di calcolo precedente. La tua proposta è molto ragionevole e include oggetti strategici. Dato che abbiamo ancora 371 punti di sforzo disponibili dal limite, propongo di aggiungere C75 che ha uno sforzo di 373, portandoci esattamente al limite massimo di 4917. Questo ci permetterebbe di massimizzare l\textbackslash{}'utilizzo del nostro budget senza sprecare capacità.'\}} \\
            }
        }
    }
    & \\ \\

    \theutterance \stepcounter{utterance}  
    & & & \multicolumn{4}{p{0.6\linewidth}}{
        \cellcolor[rgb]{0.9,0.9,0.9}{
            \makecell[{{p{\linewidth}}}]{
                \texttt{\tiny{[GM$\langle$P2]}}
                \texttt{RIFIUTO: \{'C24', 'B91', 'A15', 'A30', 'B11', 'B26', 'C47', 'A56', 'C75'\}} \\
\\ 
\texttt{PROPOSTA: \{'C24', 'B91', 'A15', 'A30', 'B11', 'B26', 'C47', 'A56', 'A19'\}} \\
\\ 
\texttt{ARGOMENTAZIONE: \{'La tua proposta supera leggermente il limite {-} il totale è 4917, che eccede di 2 punti il nostro limite di 4915. Ho sostituito C75 con A19 per rimanere entro il limite: il nuovo totale è 4907. Questa modifica mantiene tutti gli oggetti principali che abbiamo identificato come strategici e ci permette di rispettare rigorosamente i vincoli del gioco.'\}} \\
            }
        }
    }
     \\ \\

    \theutterance \stepcounter{utterance}  
    & & \multicolumn{4}{p{0.6\linewidth}}{
        \cellcolor[rgb]{0.9,0.9,0.9}{
            \makecell[{{p{\linewidth}}}]{
                \texttt{\tiny{[P1$\langle$GM]}}
                \texttt{RIFIUTO: \{'C24', 'B91', 'A15', 'A30', 'B11', 'B26', 'C47', 'A56', 'C75'\}} \\
\\ 
\texttt{PROPOSTA: \{'C24', 'B91', 'A15', 'A30', 'B11', 'B26', 'C47', 'A56', 'A19'\}} \\
\\ 
\texttt{ARGOMENTAZIONE: \{'La tua proposta supera leggermente il limite {-} il totale è 4917, che eccede di 2 punti il nostro limite di 4915. Ho sostituito C75 con A19 per rimanere entro il limite: il nuovo totale è 4907. Questa modifica mantiene tutti gli oggetti principali che abbiamo identificato come strategici e ci permette di rispettare rigorosamente i vincoli del gioco.'\}} \\
            }
        }
    }
    & \\ \\

    \theutterance \stepcounter{utterance}  
    & \multicolumn{4}{p{0.6\linewidth}}{
        \cellcolor[rgb]{0.9,0.9,0.9}{
            \makecell[{{p{\linewidth}}}]{
                \texttt{\tiny{[P1$\rangle$GM]}}
                \texttt{RIFIUTO: \{'C24', 'B91', 'A15', 'A30', 'B11', 'B26', 'C47', 'A56', 'A19'\}} \\
\\ 
\texttt{PROPOSTA: \{'C24', 'B91', 'A15', 'A30', 'B11', 'C47', 'A56', 'A19'\}} \\
\\ 
\texttt{ARGOMENTAZIONE: \{'La tua proposta supera ancora il limite {-} il totale dello sforzo è 5007, che eccede di 92 punti il nostro limite di 4915. Ho rimosso B26 dalla selezione per rispettare rigorosamente il vincolo: il nuovo totale è 4062, ben dentro il limite. Questa combinazione mantiene la maggior parte degli oggetti strategici che abbiamo identificato come prioritari, inclusi gli oggetti a basso sforzo come C24 e quelli di alto valore come A30 e B11.'\}} \\
            }
        }
    }
    & & \\ \\

    \theutterance \stepcounter{utterance}  
    & & \multicolumn{4}{p{0.6\linewidth}}{
        \cellcolor[rgb]{0.9,0.9,0.9}{
            \makecell[{{p{\linewidth}}}]{
                \texttt{\tiny{[GM$\rangle$P2]}}
                \texttt{RIFIUTO: \{'C24', 'B91', 'A15', 'A30', 'B11', 'B26', 'C47', 'A56', 'A19'\}} \\
\\ 
\texttt{PROPOSTA: \{'C24', 'B91', 'A15', 'A30', 'B11', 'C47', 'A56', 'A19'\}} \\
\\ 
\texttt{ARGOMENTAZIONE: \{'La tua proposta supera ancora il limite {-} il totale dello sforzo è 5007, che eccede di 92 punti il nostro limite di 4915. Ho rimosso B26 dalla selezione per rispettare rigorosamente il vincolo: il nuovo totale è 4062, ben dentro il limite. Questa combinazione mantiene la maggior parte degli oggetti strategici che abbiamo identificato come prioritari, inclusi gli oggetti a basso sforzo come C24 e quelli di alto valore come A30 e B11.'\}} \\
            }
        }
    }
    & \\ \\

    \theutterance \stepcounter{utterance}  
    & & & \multicolumn{4}{p{0.6\linewidth}}{
        \cellcolor[rgb]{0.9,0.9,0.9}{
            \makecell[{{p{\linewidth}}}]{
                \texttt{\tiny{[GM$\langle$P2]}}
                \texttt{ACCORDO: \{'C24', 'B91', 'A15', 'A30', 'B11', 'C47', 'A56', 'A19'\}} \\
\\ 
\texttt{ARGOMENTAZIONE: \{'Perfetto, la tua proposta rispetta il limite di sforzo e include una selezione molto equilibrata di oggetti strategici. Condividiamo chiaramente priorità simili su molti di questi elementi, e questa combinazione sembra rappresentare un buon compromesso che massimizza il valore per entrambi rimanendo ben dentro i vincoli. Accetto questa proposta.'\}} \\
            }
        }
    }
     \\ \\

    \theutterance \stepcounter{utterance}  
    & & & \multicolumn{2}{p{0.3\linewidth}}{
        \cellcolor[rgb]{0.9,0.9,0.9}{
            \makecell[{{p{\linewidth}}}]{
                \texttt{\tiny{[GM$|$GM]}}
                \texttt{game successful} \\
            }
        }
    }
    & & \\ \\

    \theutterance \stepcounter{utterance}  
    & & & \multicolumn{2}{p{0.3\linewidth}}{
        \cellcolor[rgb]{0.9,0.9,0.9}{
            \makecell[{{p{\linewidth}}}]{
                \texttt{\tiny{[GM$|$GM]}}
                \texttt{end game} \\
            }
        }
    }
    & & \\ \\

\end{supertabular}
}

\end{document}
